% Baranov, Sec. 1.4.2, Fig. 1.14. Native size: 140 x 36 mm.
% r0 = 6700 km, Delta V_t = 0.002 km/s, mu = 398600.4418 km^3/s^2.
% eta = (Delta V_t/n0)(4 sin theta - 3 theta),
% Delta r = 2(Delta V_t/n0)(1-cos theta); axes both in kilometres.
\begin{tikzpicture}[course diagram,x=1mm,y=1mm]
\path[use as bounding box] (-122,-8) rectangle (18,28);
\pgfmathsetmacro{\impulselength}{1.7372939683809365}
\draw[vector line] (-119,0)--(15,0) node[above] {$\eta$, км};
\draw[vector line] (0,-5)--(0,25) node[above left] {$\Delta r$, км};
\foreach \a in {-40,-30,-20,-10}{
 \draw[courseink,line width=.5pt] ({2.6*\a},-.8)--({2.6*\a},.8);
 \node[below=1.2mm,inner sep=0pt] at ({2.6*\a},0) {$\a$};}
\foreach \a in {3.5,7}{
 \draw[courseink,line width=.5pt] (-.8,{3*\a})--(.8,{3*\a});
 \node[right=1.2mm,inner sep=0pt] at (0,{3*\a}) {$\a$};}
\begin{scope}
\clip (-119,-.3) rectangle (6,24);
\draw[orbit line] plot[domain=0:360,samples=181,variable=\a]
 ({2.6*\impulselength*(4*sin(\a)-3*\a*pi/180)},
  {6*\impulselength*(1-cos(\a))});
\draw[courseaccent,line width=.9pt,dashed] plot[domain=360:550,samples=101,variable=\a]
 ({2.6*\impulselength*(4*sin(\a)-3*\a*pi/180)},
  {6*\impulselength*(1-cos(\a))});
\draw[courseaccent,line width=1pt,->] plot[domain=114:122,samples=9,variable=\a]
 ({2.6*\impulselength*(4*sin(\a)-3*\a*pi/180)},
  {6*\impulselength*(1-cos(\a))});
\end{scope}
\fill[courseaccent] (0,0) circle[radius=.65];
\node[below right,inner sep=1mm] at (0,0) {$0$};
\fill[courseaccent] ({-2.6*\impulselength*6*pi},0) circle[radius=.65];
\node[below=1.2mm,inner sep=0pt] at ({-2.6*\impulselength*6*pi},0) {$T_0$};
\node[text=coursemuted,font=\fontsize{8}{10}\selectfont] at (-103,25)
 {следующий виток};
\end{tikzpicture}
